%% USPSC-Abstract.tex
%\autor{Silva, M. J.}
\begin{resumo}[Abstract]
 \begin{otherlanguage*}{english}
	\begin{flushleft}
		\setlength{\absparsep}{0pt} % ajusta o espaçamento dos parágrafos do resumo
 		\SingleSpacing  		\imprimirautorabr~~\textbf{\imprimirtitleabstract}.	\imprimirdata.  \pageref{LastPage} p.
		%Substitua p. por f. quando utilizar oneside em \documentclass
		%\pageref{LastPage} f.
		\imprimirtipotrabalhoabs~-~\imprimirinstituicao, \imprimirlocal, 	\imprimirdata.
 	\end{flushleft}
	\OnehalfSpacing
   Transaction descriptions in credit card statements are short, non-standardized and noisy: they contain
   truncated abbreviations, payment intermediary prefixes and errors introduced during document reading. This
   noise prevents the automatic categorization of expenses and, with it, personal budget planning. The objective
   of this research was to design, implement and experimentally evaluate an end-to-end computational solution
   for the extraction, sanitization and automatic classification of expenses in the Brazilian context, executable
   on local infrastructure and compliant with the Brazilian General Data Protection Law. The solution combines
   document extraction with tabular layout preservation, deterministic lexical normalization routines and an
   agent orchestrated as a state graph, which articulates dense vector retrieval, reasoning with a local large
   language model and human confirmation in low-confidence cases, on a layered data infrastructure with
   experiment tracking. The evaluation used two sets of real data: 699 transactions audited against the issuing
   institution's official files, to measure extraction quality, and 1,660 transactions labeled into 11 categories, to
   compare five architectures under a merchant-disjoint blind test. Layout-preserving extraction reduced the
   character error rate from 52.82\% to 5.72\% and raised the exact match rate from 3.3\% to 86.2\%. In
   classification, the sparse representation based on character $n$-grams combined with Random Forest achieved
   the best Macro-F1 (69.3\%), with no statistically significant difference from dense vector retrieval in its best
   configuration (66.8\%; McNemar's test, $p = 0.86$), and ahead of the hybrid agent (60.0\%) and the fine-tuning
   of a byte-level model (48.3\%). The central hypothesis, stated for the hybrid pipeline, was refuted in both of
   its sub-hypotheses: no architecture reached the 85\% Macro-F1 threshold. The finding of widest
   methodological reach is the gap of 23 to 25 percentage points between conventional cross-validation and the
   merchant-disjoint blind test, which quantifies how much of the apparent performance of this class of
   classifier stems from memorizing previously observed merchants. It is concluded that the proposed pipeline
   is functionally viable and auditable under local hardware constraints, and that its predictive effectiveness
   remains conditioned on the volume and coverage of the labeled corpus.

   \vspace{\onelineskip}

   \noindent
   \textbf{Keywords}: Financial transaction classification. Natural language processing. Vector similarity search. Optical character recognition. Intelligent agents.
 \end{otherlanguage*}
\end{resumo}
